This first paper has deliberately remained close to the minimal skeleton of the theory. Its purpose has not been to complete every branch of Arithmetic Expression Geometry at once, but to isolate the few structures that appear to be unavoidable once arithmetic evaluation is treated as a geometry of histories rather than as a disposable intermediate process.

\subsection{What has been established}

The logical chain of the paper can now be read as a single progression:
\[
\text{expression paths}
\;\longrightarrow\;
\text{flow}
\;\longrightarrow\;
\text{torsion}
\;\longrightarrow\;
\text{contact structure}
\;\longrightarrow\;
\delta\text{-calculus}
\;\longrightarrow\;
\text{arithmetic holomorphicity}.
\]
Section~\ref{sec:flowequation} extracted the infinitesimal propagation law of arithmetic evaluation. Section~\ref{sec:01kind} showed that this law is not merely formal, but admits concrete geometric realization on the basic backgrounds $\mathfrak{E}_0$ and $\mathfrak{E}_1$. Section~\ref{sec:acspace} globalized the first defect of the theory by expressing torsion in the Accumulative Commutative Space as an algebraic difference, a boundary integral, and a weighted area integral. Section~\ref{sec:ecstructure} then localized this defect in the contact form
\[
\alpha = da-(\mu\,du+\lambda a\,dv),
\]
while Section~\ref{ch:differential_calculus} identified the associated horizontal differential calculus. Finally, Section~\ref{ch:holomorphic} showed that the same structure already supports a first analytic field equation and even forces a nontrivial twisted second-order consequence.

In this sense, Paper I has established more than a collection of isolated constructions. It has shown that arithmetic histories carry a coherent geometric package: a propagation law, a defect quantity, a local geometric engine, and the beginning of a function theory. That is the conceptual threshold this paper was meant to cross.

\subsection{From paths to relations}

The present paper studies \emph{paths}. The next stage should study \emph{relations}. A path records a computation history; a closed word records a relation among histories. The natural primitive object here is not equality alone, but \emph{neutrality}: the various senses in which a nontrivial history can return to a null effect.

At the most elementary level, four kinds of neutrality already suggest themselves.

\begin{enumerate}
    \item \textbf{Value neutrality.} A closed history returns the same final assignment value.
    \item \textbf{Operator neutrality.} The induced transformation on assignment values is the identity.
    \item \textbf{Charge neutrality.} The total additive and multiplicative charges vanish in the commutative record.
    \item \textbf{Torsion neutrality.} The enclosed weighted area, hence the arithmetic torsion, vanishes.
\end{enumerate}

These notions should not be collapsed into one another. Two histories may agree in final value while remaining distinct in charge or torsion, and it is precisely such cases that are likely to produce the first serious theory of \emph{meaningful zeros}. Put differently, equality asks when two expressions are identified as the same object, whereas loop theory asks in what sense a nontrivial relation disappears and what residue of that disappearance remains visible.

This is why the loop direction should be viewed as a relation theory rather than merely as a supplement to the present path language. Once a corpus of explicit closed words is assembled, one can begin to classify which neutralities imply which others, which implications fail, and which loops remain nontrivial even when one layer of effect has vanished. That program is the natural algebraic and topological continuation of the current paper.

\subsection{From differential geometry to function theory}

The second major direction opened here is analytic rather than relational. The operator package
\[
\bar\partial_{\mathrm{AEG}}F=0,
\qquad
4\partial_{\mathrm{AEG}}\bar\partial_{\mathrm{AEG}}
=
\Delta_H+i\mu\lambda\,\partial_a
\]
shows that arithmetic holomorphicity is not decorative notation. It is the first sign that AEG may support a genuine function theory.

Several near-term analytic questions are now concrete enough to be pursued directly. One must determine the canonical second-order operators attached to the arithmetic horizontal structure; construct explicit families of arithmetic holomorphic and twisted harmonic fields; clarify the role of harmonic conjugacy and local normal forms; and eventually formulate the first kernel, boundary value, or continuation problems on the basic backgrounds. The affine--Appell basis is especially promising here, because it gives the theory a computational engine before it has a large abstract apparatus.

The main discipline of this analytic route is equally clear. It should not diffuse into a general theory of contact or sub-Riemannian analysis detached from arithmetic meaning. The constants $\mu$ and $\lambda$, the value-dependent scaling direction, and the interpretation of $D_u$ and $D_v$ as add- and multiply-directions must remain visible at every stage. Otherwise the theory will become formally rich but conceptually misplaced.

\subsection{Singular backgrounds and higher structures}

The first-kind space $\mathfrak{E}_1$ already suggests that singular backgrounds are not decorative variants of $\mathfrak{E}_0$. They are likely to be the place where torsion, neutrality, and analysis interact most sharply. A loop may be trivial at one level and obstructed at another because of singular geometry; an analytic field may extend smoothly on a regular background but fail, branch, or pick up monodromy around a singular locus. These possibilities do not yet constitute a theory, but they indicate where the next structural collisions are likely to occur.

Beyond loops and singular backgrounds lies the possibility of arithmetic $2$-cells and arithmetic surfaces. If such objects become viable, they may provide preferred fillings for loops, new area-type invariants, and a higher-dimensional stage on which zero classes, holonomy, and analytic structures can meet. For the present paper, however, it is enough to note that the basic ingredients are already visible: boundary histories, weighted area, contact curvature, and a first horizontal complex calculus.

\subsection{A realistic program after Paper I}

A decisive next stage would already be achieved if the following four goals were met.

\begin{enumerate}
    \item Build a reliable corpus of explicit closed words and separate the first neutrality classes by concrete examples.
    \item Fix a canonical first-order/second-order analytic package on the basic backgrounds and construct nontrivial explicit solution families.
    \item Prove at least one theorem connecting the relation theory back to the analytic or geometric theory, for example by attaching an invariant to a loop or by constraining fillings through torsion or singularity data.
    \item Clarify the role of singular backgrounds as genuine structural objects rather than as exceptional cases.
\end{enumerate}

If these goals are reached, Arithmetic Expression Geometry will have moved beyond a geometric reinterpretation of evaluation and toward a unified framework in which arithmetic histories admit propagation, defect, relation, and analysis in a single language.

\medskip

The broader claim of this paper is therefore modest in form but ambitious in implication. Arithmetic expressions do not merely yield values; they also carry histories, and those histories have geometry. Once that geometry is taken seriously, arithmetic begins to exhibit flow, torsion, horizontal structure, curvature, and analytic fields. This paper has only established the first layer of that picture. But if the subsequent relation theory and analytic theory develop as expected, then the geometry of arithmetic expressions may become not just an analogy, but a stable mathematical domain of its own.
